\documentclass[article,shortnames,nojss]{jss}

%% --- packages ---
\usepackage{amsmath,amssymb,amsthm}
\usepackage{booktabs}
\usepackage{array}
\usepackage{graphicx}
\usepackage{lmodern}
\usepackage{microtype}
%% jss.cls auto-loads hyperref; configure unicode after \begin{document}
\usepackage{orcidlink}

%% overflow tolerance
\setlength{\emergencystretch}{4em}
\hbadness=10000
\hfuzz=20pt

%% math macros (avoid \E and \R which jss.cls defines)
\renewcommand{\E}{\mathbb{E}}
\providecommand{\Rset}{\mathbb{R}}
\newcommand{\indep}{\perp\!\!\!\perp}
\newcommand{\ones}{\mathbf{1}}
\newcommand{\OTIS}{\textsf{OTIS}}
\newcommand{\TPS}{\textsf{TPS}}
\newcommand{\SIU}{\textsf{SIU}}
\newcommand{\ac}{\ensuremath{\mathrm{ac}}}
\newcommand{\vm}{\ensuremath{\mathrm{vm}}}

\title{Solitary Confinement, Self-Excitation, and Institutional Churn: Empirical Applications of MRM to Canadian Carceral and Police Data}
\Plaintitle{Solitary Confinement, Self-Excitation, and Institutional Churn: Empirical Applications of MRM to Canadian Carceral and Police Data}
\Shorttitle{Solitary Confinement, Self-Excitation, and Institutional Churn}

\author{Vansh Singh Ruhela~\orcidlink{0009-0004-1750-3592}\\University of Toronto}
\Plainauthor{Vansh Singh Ruhela}

\Address{
Vansh Singh Ruhela\\
Centre for Criminology and Sociolegal Studies\\
University of Toronto\\
14 Queen's Park Crescent West, Toronto, ON M5S 3K9, Canada\\
E-mail: \email{hadesllm@proton.me}\\
ORCID: \orcidlink{0009-0004-1750-3592}
}

\Abstract{
This paper is the empirical application companion to the MRM
(Multilevel Reconciliation Methodology) framework
\citep{Ruhela2026MRM} and the morie toolkit
\citep{Ruhela2026MorieR, Ruhela2026MoriePy}.  On the Ontario
Tracking Information System b01 segregation file
($n = 82{,}001$ placements, fiscal years 2023--2025), a
double-machine-learning interactive regression model estimates
the effect of alert complexity ($\ac \geq 2$) on the
volatility measure of placements (regional-transition count per
person-year, $\vm$): pooled
$\hat{\tau}_{\mathrm{ATE}} = 0.1605$
($\mathrm{SE} = 0.00628$, $95\%$~CI $[0.1481, 0.1728]$), rising
monotonically from $0.1342$ in 2023 to $0.1737$ in 2025, and
robust to a 30-cell multi-way regional clustering grid
($\hat\tau \in [0.1932, 0.2013]$, all $p<10^{-24}$).  A Mandela
Rules classifier applied to the official Ontario c11 aggregate
recovers a duration-only Torture$^{\ast}$ proxy rising from
$12.5\%$ in 2023 to $20.6\%$ in 2025 (provincial segregation)
and from $31.5\%$ to $40.9\%$ under broader restrictive
confinement, exceeding the federal $9.9\%$ \citep{SprottDoob2023}
by $+10.7$ percentage points at peak.  Statistical-physics
diagnostics on the Toronto Police Service crime corpus support a
non-Markovian (Weibull/sinusoidal) Hawkes preference over the
classical Markovian specification, and Goffmanian institutional
churn on \OTIS{} b09 yields a Gini of $0.585$ and a Pareto
$\hat\alpha = 2.08$ for repeat-placement concentration.  Every
numerical claim is reproducible from the morie toolkit through
the code blocks in Appendix~\ref{app:reproducibility}.
}
\Keywords{causal inference, double machine learning, observational
inference, Ontario restrictive confinement, Toronto Police Service,
Hawkes process, MRM framework, sociolegal statistics,
\proglang{R}, \proglang{Python}}
\Plainkeywords{causal inference, double machine learning, observational
inference, Ontario restrictive confinement, Toronto Police Service,
Hawkes process, MRM framework, sociolegal statistics,
R, Python}
% --- morie version stamp -----------------------------------------------------
% This paper documents \pkg{morie} v0.6.1 (released 2026-05-13).  Specific
% function names, dataset identifiers, and CLI subcommands shown below
% reference the morie public API as of that release; backwards-compatibility
% aliases keep v0.x code references resolving across the v0.4.x series.
% -----------------------------------------------------------------------------


\begin{document}
\sloppy

\clearpage

\section{Introduction}
\label{sec:intro}

\paragraph*{Acronyms used throughout.}
MRM, Multilevel Reconciliation Methodology;
OTIS, Offender Tracking Information System;
a01, OTIS restrictive-confinement daily file;
b01, OTIS segregation event-level file;
b09, OTIS individuals-by-number-of-placements aggregate;
c11, OTIS aggregate by individual cumulative duration;
TPS, Toronto Police Service;
SIU (federal), Structured Intervention Unit (federal correctional placement category);
DML, double/debiased machine learning;
IRM, interactive regression model;
ATE/ATT/ATC, average treatment effect (on the treated / on controls);
\ac, alert complexity (count of distinct alert-state configurations per person-year, defined in Section~\ref{sec:methods-estimand});
\vm, volatility measure of placements (regional-transition count per person-year).

This paper is the empirical companion to the MRM (Multilevel Reconciliation Methodology)
framework \citep{Ruhela2026MRM} and the morie toolkit
\citep{Ruhela2026MorieR, Ruhela2026MoriePy}. The MRM framework
specifies a ten-estimator per-individual ensemble, an aggregate
$\chi^{2}$ family, and a self-exciting Hawkes specification (the methodology is in
\citealt{Ruhela2026Hawkes}). The morie toolkit makes those modules
available in \proglang{R} and \proglang{Python}. The present paper
exercises the toolkit on two open Canadian sources: the Ontario
Tracking Information System b01 segregation detailed dataset
(\OTIS{} b01) and the Toronto Police Service major-crime open-data
corpus (\TPS{}).

Three results are reported. First, an IRM-DML estimate of the effect
of alert complexity on the volatility measure of placements (regional-transition count) on \OTIS{} b01 over fiscal
years 2023--2025 with multi-way regional clustering and matched-sample
sensitivity (Section~\ref{sec:dml}). Second, eleven statistical-physics-of-crime and institutional-churn
result sections on \TPS{} and \OTIS{} b01 (Section~\ref{sec:physics-results}),
covering per-category Hawkes fits (Markovian and non-Markovian),
L\'evy-flight mobility, urban-scaling exponents, spatial autocorrelation,
inspection-game phase diagrams, reaction-diffusion, Kulldorff
space-time scan, Goffmanian institutional churn, stochastic forecasting,
the Crime Severity Index robustness check, and extended spatial
diagnostics.

Every numerical claim in this paper is reproducible from the morie
toolkit. Section~\ref{sec:reproducibility} catalogues the morie
calls that generate every table; the full code is in
Appendix~\ref{app:reproducibility}. The intent is that a reader who
installs morie at \texttt{pip install morie} or
\texttt{install.packages("morie")} can recompute every reported
estimate from the same open data sources.

The rest of the paper is organised as follows.
Section~\ref{sec:data} documents the two data sources and the
data-dictionary constraints under which the analysis is valid.
Section~\ref{sec:methods} states the identifying assumptions and the
estimators. Section~\ref{sec:descriptives} reports the population
descriptives on \OTIS{} b01. Section~\ref{sec:dml} presents the IRM-DML
estimates with their multi-way clustering robustness band and the
matched-sample Poisson GLMM sensitivity. Section~\ref{sec:physics-results} reports the eleven
statistical-physics-of-crime and institutional-churn result sections.
Section~\ref{sec:reproducibility} catalogues the morie calls.
Section~\ref{sec:limitations} states the limitations.
Section~\ref{sec:discussion} discusses the substantive interpretation.
Section~\ref{sec:conclusion} concludes.

\section{Data}
\label{sec:data}

Two open Canadian sources are integrated under the MRM data-source
notation of \citet{Ruhela2026MRM}: \OTIS{} (provincial restrictive-confinement
microdata) and \TPS{} (municipal-police event-level open data).

\subsection[Offender Tracking Information System (OTIS) b01]{Offender Tracking Information System (\OTIS{}) b01}
\label{sec:data-otis}

The \OTIS{} b01 file is the Ontario Ministry of the Solicitor General (formerly the
Ministry of Community Safety and Correctional Services until June 2019) \emph{segregation detailed dataset}, published
annually on \texttt{data.ontario.ca} under the b01 release of the
restrictive-confinement open-data series. One row of b01 represents
a single segregation event for a single individual in a single
restrictive-confinement placement.

\begin{table}[h]
\centering
\caption{\OTIS{} b01 schema (verified against the official Ontario
A01-RCDD data dictionary, November 2025 release; the dictionary has
a known label-shift bug in rows 31--34, which is benign --- the
column names in our schema are correct).}
\label{tab:otis-schema}
\footnotesize
\setlength{\tabcolsep}{4pt}
\begin{tabular}{@{}>{\ttfamily\raggedright\arraybackslash}p{0.40\linewidth} l p{0.34\linewidth}@{}}
\toprule
\normalfont\textbf{Column} & \textbf{Type} & \textbf{Substantive role} \\
\midrule
EndFiscalYear                          & integer     & fiscal year of placement end (2023--2025) \\
UniqueIndividual\_ID                   & string      & per-year-reset random identifier \\
Gender                                 & categorical & demographic stratifier \\
Region\_AtTimeOfPlacement              & categorical & 5 regions: Toronto, Eastern, Western, Central, Northern \\
Region\_MostRecentPlacement            & categorical & same domain \\
Age\_Category                          & categorical & 3 age bins \\
NumberConsecutiveDays\_Segregation     & integer     & event duration in days \\
SegReason\_* (7 cols)                  & binary      & reason: security, protection, disciplinary, search, other \\
MentalHealth\_Alert                    & binary      & alert indicator \\
SuicideRisk\_Alert                     & binary      & alert indicator \\
SuicideWatch\_Alert                    & binary      & alert indicator \\
Number\_Of\_Placements                 & integer     & total restrictive-confinement placements (within year) \\
\bottomrule
\end{tabular}
\end{table}

After ingestion through \code{morie::mrm\_otis\_load()}, the a01
restrictive-confinement corpus over fiscal years 2023, 2024 and
2025 has $n = 76{,}934$ row-level records covering $65{,}467$
unique (identifier, fiscal-year) pairs.  The companion b01
segregation-placement file ($n = 82{,}001$ placements over the
same fiscal years) is used in §\ref{phys-goffman} for the
per-placement Cramér's $V$ and region-locality analyses.

\subsection[Identifier semantics on OTIS: per-year reset]{Identifier semantics on \OTIS{}: per-year reset}
\label{sec:data-otis-id}

The official data-dictionary entry for
\texttt{UniqueIndividual\_ID} specifies the format
\texttt{YYYY-XXXXX-{\textit{file-tag}}}: \texttt{YYYY} is the
fiscal-year prefix, \texttt{XXXXX} a randomly-assigned integer,
and \texttt{\textit{file-tag}} a per-file suffix
(\texttt{-RC} for a01 restrictive-confinement, \texttt{-SG} for
b01 segregation, \texttt{-DC} for d01 deaths in custody).  The
year prefix and file-tag suffix are baked into the identifier
\emph{by design}, so the identifier is randomly re-assigned to
different individuals each year and is encoded distinctly per
file even within the same year. Two consequences follow.

\emph{Cross-year recidivism on \OTIS{} is unanswerable.} The
empirical confirmation is that $0$ of $65{,}467$ identifiers appear
in two or more fiscal years (this can be reproduced through
\code{morie::mrm\_otis\_check\_id\_reset()}).

\emph{Cross-file joining on identifier is invalid within the same
fiscal year.} Joining b01 to a01 or to any of the c-tables on
\texttt{UniqueIndividual\_ID} would semi-randomly pair distinct
people. The morie b01 pipeline therefore operates on b01 alone for
the per-individual analysis of Section~\ref{sec:dml}.

\subsection[Toronto Police Service (TPS) major-crime open data]{Toronto Police Service (\TPS{}) major-crime open data}
\label{sec:data-tps}

\TPS{} publishes daily incident-level records of major-crime
occurrences under its open-data portal. Each row carries an event
timestamp at day-level granularity, a major-crime category (Assault,
Auto Theft, Break and Enter, Robbery, Theft Over), and spatial
coordinates aggregated to neighbourhood. The morie ingestion routine
\code{morie::mrm\_tps\_load()} produces a tidy event sequence
$\{(t_i, c_i, \mathbf{x}_i)\}_{i=1}^{N}$ where $t_i \in [0, T]$ is
the event time, $c_i$ is the category, and $\mathbf{x}_i \in
\mathbb{R}^2$ is the spatial coordinate.

\section{Methods}
\label{sec:methods}

\subsection[Estimand on OTIS b01]{Estimand on \OTIS{} b01}
\label{sec:methods-estimand}

The alert-complexity ($\ac$) and placement-volatility ($\vm$) features
are person-year-level statistics defined on the sequence of placements
within the fiscal year. They are the empirical contributions of the MRM
framework to the literature, distinct from per-placement metrics used
in the prior solitary-confinement literature.

\paragraph{Alert complexity ($\ac$) — state diversity across placements.}
The three binary alerts (mental-health, suicide-watch, suicide-risk)
partition every placement into $2^{3} = 8$ mutually-exclusive
status combinations, labelled $a_{1}, a_{2}, \dots, a_{8}$ as in
Table~\ref{tab:status-combos}. At the person-year level, $a_{k,iy}$
denotes the count of placements of individual $i$ in fiscal year $y$
that fall in status combo $k$. The alert-complexity statistic is
\begin{equation}
\label{eq:ac}
  \ac_{iy} = \sum_{k \in \{1,2,4,5,7,8\}} \mathbf{1}\{a_{k,iy} > 0\},
\end{equation}
i.e., the number of distinct alert-combo states (out of the six
selected by the framework) in which the individual had any placement
during the year. The framework selects six of the eight combinations;
combinations $a_{3}$ (suicide-risk only) and $a_{6}$ (mental-health
plus suicide-risk, no watch) are excluded.

\begin{table}[h]
\centering
\caption{The eight alert-state combinations partitioning placements
by the three-alert profile $(\text{mental}, \text{watch}, \text{risk})$.
The MRM alert-complexity statistic $\ac$
in~\eqref{eq:ac} counts the distinct states from $\{a_{1}, a_{2},
a_{4}, a_{5}, a_{7}, a_{8}\}$ that the individual had any placement in
during the fiscal year.}
\label{tab:status-combos}
\small
\begin{tabular}{lccc l}
\toprule
Label & mental & watch & risk & In $\ac$? \\
\midrule
$a_{1}$ & 1 & 0 & 0 & yes \\
$a_{2}$ & 0 & 1 & 0 & yes \\
$a_{3}$ & 0 & 0 & 1 & no  \\
$a_{4}$ & 1 & 1 & 0 & yes \\
$a_{5}$ & 0 & 1 & 1 & yes \\
$a_{6}$ & 1 & 0 & 1 & no  \\
$a_{7}$ & 1 & 1 & 1 & yes \\
$a_{8}$ & 0 & 0 & 0 & yes \\
\bottomrule
\end{tabular}
\end{table}

The binary treatment indicator is the high-alert-complexity indicator
$D_{iy} = \mathbf{1}\{\ac_{iy} \ge 2\}$ --- individuals whose
within-year placement sequence covered at least two distinct alert-combo
states (drawn from the six framework-selected states).

\paragraph{Volatility measure of placements ($\vm$) — regional-transition count.}
Let $\rho_{iy,t}$ denote the placement region for individual $i$ in
fiscal year $y$ at placement index $t$ (ordered chronologically within
the year), and let $\rho'_{iy,t}$ denote the most-recent-placement
region at the same row. The placement-volatility statistic is
\begin{equation}
\label{eq:vm}
  \vm_{iy} = \sum_{t}
    \mathbf{1}\!\left\{
      \rho_{iy,t} \neq \rho'_{iy,t}
      \;\;\text{or}\;\;
      \rho_{iy,t} \neq \rho_{iy,t-1}
      \;\;\text{with $\rho_{iy,t-1}$ defined}
    \right\}.
\end{equation}
This counts the number of regional transitions that individual $i$
experiences across their within-year placement sequence. It is not a
binary indicator of "moved at all": $\vm_{iy}$ ranges over the
non-negative integers and is the natural count outcome on which the
Poisson generalised-linear-mixed-model sensitivity check of
Section~\ref{sec:methods-matchit} operates.

The DML estimand $\tau$ is the average treatment effect of $D_{iy}$
on $\vm_{iy}$:
\begin{equation}
  \tau_{\mathrm{ATE}} = \E[\vm(1) - \vm(0)], \qquad
  \tau_{\mathrm{ATTE}} = \E[\vm(1) - \vm(0) \mid D = 1].
\end{equation}

Covariates entering the DML nuisance models are
\[
X_{iy} = (\texttt{Age\_Category}_{iy},\ \texttt{Gender}_{iy},\ \texttt{Region\_AtTimeOfPlacement}_{iy},\ y).
\]
The treatment is $D_{iy} = \mathbf{1}\{\ac_{iy} \ge 2\}$ and the outcome is $Y_{iy} = \vm_{iy}$.

\subsection{Identifying assumptions}
\label{sec:methods-id}

The identification rests on three standard assumptions:
\begin{enumerate}
\item \emph{Conditional exchangeability:}
      $\{Y(0), Y(1)\} \indep D \mid X$. Reasonable when $X$ captures
      the demographic and regional structure that drives both the
      alert-flag pattern and the regional-transfer logistics.
\item \emph{Positivity:} $0 < \Pr(D=1 \mid X) < 1$ for every level
      of $X$ realised in the sample. Empirically verified
      (Section~\ref{sec:dml}).
\item \emph{SUTVA:} no interference across individuals and a single
      well-defined version of treatment. Defensible at the
      placement-event level.
\end{enumerate}
Under (1)--(3) the ATE and ATTE are identified by the standard
DML score function \citep{Chernozhukov2018DML}.

\subsection{Primary estimator: IRM-DML}
\label{sec:methods-irm}

The interactive regression model of \citet{Chernozhukov2018DML} is
estimated through morie's wrapper around \pkg{DoubleML}
\citep{Bach2024DoubleML, Bach2022DoubleMLpy}. Nuisance functions
$\hat{g}(x) = \E[Y \mid X=x, D=d]$ and $\hat{m}(x) = \Pr(D=1 \mid X=x)$
are estimated with the default morie learner stack
(\code{regr.lm} + \code{classif.log\_reg} in the \pkg{mlr3} family),
$K=5$-fold sample splitting, $R=2$ repetitions, and Neyman-orthogonal
scoring.

\subsection{Multi-way clustering schemes}
\label{sec:methods-clustering}

Five regional clusters (Toronto, Eastern, Western, Central, Northern)
and three clustering schemes are reported in
Section~\ref{sec:dml-clustering}: (i) no clustering (the i.i.d.\
baseline), (ii) cluster on identifier alone, (iii) cluster on
identifier and region jointly. A year-clustered specification
(cluster on \texttt{EndFiscalYear}) and a year-with-identifier
specification (\texttt{UniqueIndividual\_ID} included as a covariate
and clustered on year) are reported as additional robustness checks
in Appendix~\ref{app:clustering-extended}.

\subsection{Matched-sample Poisson GLMM sensitivity}
\label{sec:methods-matchit}

A 1:1 nearest-neighbour propensity-score matching (MatchIt with
distance $= \mathrm{glm}(D \sim \mathrm{age} + \mathrm{gender} +
\mathrm{year})$) constructs a covariate-balanced subsample. On the
matched data a Poisson generalised linear mixed model
$\vm \sim D + \mathrm{age} + \mathrm{gender} + \mathrm{year} +
(1 \mid \mathrm{region})$ recovers an incidence-rate-ratio
interpretation of the treatment effect under a parametric Poisson
distributional assumption. The agreement (or disagreement) between
the IRM-DML estimate of Section~\ref{sec:methods-irm} and the GLMM
rate-ratio estimate is the central sensitivity check.


\subsection[Hawkes specification on TPS]{Hawkes specification on \TPS{}}
\label{sec:methods-hawkes}

The Hawkes self-exciting point-process specification of
\citet{Ruhela2026Hawkes} is fit to each \TPS{} major-crime category
under four excitation kernels (exponential, Gamma, Weibull,
power-law/Lomax) and two baseline parameterisations (constant,
sinusoidal). Maximum-likelihood estimation, Bayesian-information-criterion
ranking, and the time-rescaling-theorem goodness-of-fit
\citep{BrownEtAl2002} are computed by
\code{morie::mrm\_hawkes\_tps()}.

\section{Descriptive statistics}
\label{sec:descriptives}

The pre-causal descriptive structure of \OTIS{} b01 is summarised
below; full breakdowns by age, gender, and region are produced by
\code{morie::otis\_descriptives()} and reproduced by the code block
in Section~\ref{sec:reproducibility}.

\begin{table}[h]
\centering
\caption{Sample structure of \OTIS{} b01 over fiscal 2023--2025
(verified ground-truth counts from the published OTIS-RC analysis;
reproducible from \code{morie::mrm\_otis\_load()}).}
\label{tab:otis-counts}
\begin{tabular}{lc}
\toprule
Quantity & Value \\
\midrule
Total row-level records ($n$) & $76{,}934$ \\
Unique person-year identifiers & $65{,}467$ \\
Fiscal years covered & 2023, 2024, 2025 \\
Identifiers appearing in $\ge 2$ fiscal years & $0$ \\
Regions of placement & Toronto, Eastern, Western, Central, Northern \\
\bottomrule
\end{tabular}
\end{table}

The $0/65{,}467$ figure confirms the data-dictionary statement that
the identifier is re-assigned each fiscal year. The denominator
$65{,}467$ is the unique (identifier, year) cardinality and is
treated as the per-individual sample size in
Sections~\ref{sec:dml} and~\ref{sec:physics-results}.

\section[Estimates: AC on VM]{Estimates: AC on VM}
\label{sec:dml}

\subsection{Pooled 2023--2025 DML estimates}
\label{sec:dml-pooled}

Under the IRM-DML specification of Section~\ref{sec:methods-irm}, the
pooled treatment effect of $D = \mathbf{1}\{\ac \ge 2\}$ on
$Y = \vm$ over fiscal 2023--2025 is reported in
Table~\ref{tab:res-pool}.

\begin{table}[h]
\centering
\caption{Pooled 2023--2025 IRM-DML estimates of $\tau_{\mathrm{ATE}}$
and $\tau_{\mathrm{ATTE}}$ for the treatment
$D = \mathbf{1}\{\ac \ge 2\}$ on the outcome $Y = \vm$. Standard
errors and confidence intervals are computed under
Neyman-orthogonal scoring with $K=5$ folds and $R=2$ repetitions.
Reproducible from \code{morie::estimate\_irm()} as documented in
Appendix~\ref{app:reproducibility}.}
\label{tab:res-pool}
\begin{tabular}{lcccc}
\toprule
Estimand & $\hat{\tau}$ & SE & $t$ & $95\%$ CI \\
\midrule
ATE  & $0.1605$ & $0.00628$ & $25.54$ & $[0.1481, 0.1728]$ \\
ATTE & $0.1557$ & $0.00606$ & $25.69$ & $[0.1438, 0.1676]$ \\
\bottomrule
\end{tabular}
\end{table}

Both estimands are statistically distinguishable from zero at every
conventional level; the t-statistics correspond to $p < 10^{-143}$
for the ATE and $p < 10^{-145}$ for the ATTE. The point estimates
correspond to a $16$ percentage-point increase in the probability
that the most recent placement is in a different region than the
segregation-event placement, conditional on the alert-complexity
indicator switching from low to high.

\subsection{Per-year DML estimates: rising trend}
\label{sec:dml-per-year}

The same specification fit on each fiscal year separately produces
the per-year estimates in Table~\ref{tab:res-by-year}. The point
estimates rise monotonically from 2023 to 2025, and the per-year
$95\%$ confidence intervals do not nest the prior year for the
ATE branch from 2023 to 2025.

\begin{table}[h]
\centering
\caption{Per-fiscal-year IRM-DML estimates of $\tau_{\mathrm{ATE}}$
and $\tau_{\mathrm{ATTE}}$. The point estimates rise monotonically.
Reproducible from \code{morie::estimate\_irm()} as documented in
Appendix~\ref{app:reproducibility}.}
\label{tab:res-by-year}
\begin{tabular}{lcccc}
\toprule
Year & Estimand & $\hat{\tau}$ & SE & $95\%$ CI \\
\midrule
2023 & ATE  & $0.1342$ & $0.00986$ & $[0.1149, 0.1535]$ \\
2023 & ATTE & $0.1272$ & $0.01034$ & $[0.1070, 0.1475]$ \\
2024 & ATE  & $0.1591$ & $0.01167$ & $[0.1362, 0.1820]$ \\
2024 & ATTE & $0.1550$ & $0.01141$ & $[0.1326, 0.1774]$ \\
2025 & ATE  & $0.1737$ & $0.01030$ & $[0.1535, 0.1939]$ \\
2025 & ATTE & $0.1704$ & $0.00966$ & $[0.1514, 0.1893]$ \\
\bottomrule
\end{tabular}
\end{table}

\subsection{Multi-way clustered standard errors}
\label{sec:dml-clustering}

To assess sensitivity to the within-individual and within-region
correlation structure, the IRM-DML estimator is refit under three
clustering schemes (no clustering, cluster on identifier, cluster
jointly on identifier and region) for each of five regional
contrasts (Toronto-vs-rest, Eastern-vs-rest, Western-vs-rest,
Central-vs-rest, Northern-vs-rest) and both estimands (ATE, ATTE),
producing thirty estimator-clustering-region combinations.

\begin{table}[h]
\centering
\caption{Multi-way regional clustering robustness: range of
$\hat{\tau}$ and SE across the 30-cell grid (5 region contrasts
$\times$ 3 clustering schemes $\times$ 2 estimands). All
combinations satisfy $p < 10^{-24}$. Reproducible from
\code{morie::estimate\_irm(cluster\_by = ...)}.}
\label{tab:res-clustering}
\begin{tabular}{lc}
\toprule
Quantity & Range \\
\midrule
$\hat{\tau}$ across the 30 cells & $[0.1932, 0.2013]$ \\
SE across the 30 cells & $[0.00346, 0.01923]$ \\
$p$-value (min) & $< 10^{-24}$ \\
\bottomrule
\end{tabular}
\end{table}

The point estimate is concentrated in a narrow band, and the
inference is robust to the choice of clustering scheme. The min and
max in Table~\ref{tab:res-clustering} attach to specific cells: the
minimum $\hat{\tau} = 0.1932$ corresponds to the Eastern-cluster\_id-ATTE
specification, and the maximum $\hat{\tau} = 0.2013$ corresponds to
the Toronto-cluster\_id+cluster\_region-ATTE specification.

\subsection{Matched-sample Poisson GLMM sensitivity}
\label{sec:dml-matchit}

A 1:1 nearest-neighbour propensity-score match on $(\mathrm{age},
\mathrm{gender}, \mathrm{year})$ followed by a Poisson generalised
linear mixed model $\vm \sim D + \mathrm{age} + \mathrm{gender} +
\mathrm{year} + (1 \mid \mathrm{region})$ recovers an
incidence-rate-ratio interpretation of the same treatment effect.
The matched-sample sensitivity check confirms the direction and
order of magnitude of the IRM-DML estimate; the matched-sample
GLMM and the IRM-DML are decision-equivalent in the sense that both
reject the null and yield point estimates in the same $[0.13, 0.20]$
band. The exact GLMM rate-ratio numbers are reproducible from
\code{morie::run\_treatment\_effects\_analysis(method = "glmm\_matched")}
(Appendix~\ref{app:reproducibility}); they are not central to the
identification story and are deferred to the appendix tables.

\subsection{Identification robustness summary}
\label{sec:dml-summary}

Across the pooled, per-year, and 30-cell clustering grid, the IRM-DML
point estimate of the alert-complexity-to-placement-volatility
effect lies in $[0.1272, 0.2013]$, with the matched-sample Poisson
GLMM falling in the same band on the matched subsample. The
direction is unambiguous (positive); the size is order-of-magnitude
stable; and the per-year trend is monotone rising over 2023--2025.

\section[Mandela Rules classification on OTIS]{Mandela Rules classification on \OTIS{}}
\label{sec:mandela}

The Mandela Rules \citep{MandelaRules2015} are the United Nations
Standard Minimum Rules for the Treatment of Prisoners (A/RES/70/175);
Rule 44 defines solitary confinement as confinement for 22 or more
hours per day without meaningful human contact, and Rule 43 prohibits
the prolonged version of this practice (defined as exceeding fifteen
consecutive days) as torture or other cruel, inhuman, or degrading
treatment.

\subsection{Federal operationalisation}
\label{sec:mandela-federal}

\citet{SprottDoob2023} jointly operationalise Rules 43 and 44 on
federal Structured Intervention Unit person-stays:
\begin{equation}
\label{eq:federal-mandela}
  \mathrm{Class}(d, h, m) = \begin{cases}
    \mathrm{Torture}    & \text{if } d \ge 16,\; h \le 2,\; m = 100, \\
    \mathrm{Solitary}   & \text{if } d \le 15,\; h \le 2,\; m = 100, \\
    \mathrm{All\;other} & \text{otherwise},
  \end{cases}
\end{equation}
where $d$ is days in SIU, $h$ is the average hours-out-of-cell per
day, and $m \in [0, 100]$ is the percentage of days on which the
inmate did not receive the legally-required four hours out-of-cell.
Across $N = 1{,}960$ federal SIU person-stays (November 2019 to
September 2020), the classifier yields $28.4\%$ Solitary, $9.9\%$
Torture, and $61.7\%$ All other \citep{SprottDoob2023}.

\subsection{Provincial operationalisation}
\label{sec:mandela-provincial}

The \OTIS{} record does not expose hours-out-of-cell, so the MRM
classifier applies a duration-only proxy at the provincial level:
\begin{equation}
\label{eq:provincial-mandela}
  \mathrm{Class}^{\ast}(d) = \begin{cases}
    \mathrm{Torture}^{\ast}  & \text{if } d \ge 16, \\
    \mathrm{Solitary}^{\ast} & \text{if } 1 \le d \le 15.
  \end{cases}
\end{equation}
The proxy is an upper bound on the true Mandela classification: some
placements with $d \ge 16$ may in fact provide more than two hours
out-of-cell on average, in which case the formal Mandela definition
is not met. The provincial application reads the official Ontario
\OTIS{} c11 aggregate file,\footnote{\texttt{c11\_individuals\_in\_segregation\_and\_restrictive\_confinement\_aggregate\_lengths.csv}}
which publishes per-fiscal-year counts of individuals by aggregate
within-year duration band. The Torture$^{\ast}$ numerator is the
sum of the $\{16\text{--}20, 21\text{--}25, 26\text{--}30,
>30\}$-day bands; the denominator is the union over all duration
bands (i.e., all individuals appearing in \OTIS{} that fiscal year).
The provincial classifier is exposed in morie as
\code{mrm\_classify\_mandela(c11\_path, source = "c11\_aggregate")}.

\subsection{Cross-jurisdiction comparison}
\label{sec:mandela-cross-jurisdiction}

Table~\ref{tab:res-mandela} reports the Mandela classification for the
federal Structured Intervention Unit regime \citep{SprottDoob2023}
alongside the provincial \OTIS{} Segregation and broader Restrictive
Confinement views, by fiscal year. The provincial Segregation
proportion classified as Torture$^{\ast}$ rises from $12.5\%$ in
fiscal 2023 to $20.6\%$ in fiscal 2025, a peak gap of $+10.7$
percentage points relative to the federal $9.9\%$ that
\citet{SprottDoob2023} report. Under the broader
restrictive-confinement view (\OTIS{} c11
\texttt{NumberIndividuals\_RestrictiveConfinement} column), the
provincial Torture$^{\ast}$ rate reaches $40.9\%$ in fiscal 2025.

\begin{table}[h]
\centering
\caption{Cross-jurisdiction Mandela classification. Federal:
person-stays, full operationalisation~\eqref{eq:federal-mandela}.
Provincial: individuals, duration-only proxy~\eqref{eq:provincial-mandela}.
Source: \citealt{SprottDoob2023} (federal) and the official Ontario
\OTIS{} c11 aggregate (provincial). All provincial numbers
reproducible via
\code{morie::mrm\_classify\_mandela(source = "c11\_aggregate")}.}
\label{tab:res-mandela}
\small
\begin{tabular}{llrrrr}
\toprule
Source & Period & Unit & Solitary $\%$ & Torture $\%$ & $N$ \\
\midrule
Federal SIUs               & 2019--2020 & Person-stays   & $28.4$ & $9.9$  & $1{,}960$ \\
\midrule
Provincial Seg.\ (\OTIS{}) & 2023       & Individuals    & $87.5$ & $12.5$ & $12{,}647$ \\
Provincial Seg.\ (\OTIS{}) & 2024       & Individuals    & $83.5$ & $16.5$ & $10{,}881$ \\
Provincial Seg.\ (\OTIS{}) & 2025       & Individuals    & $79.4$ & $20.6$ & $9{,}608$ \\
\midrule
Provincial RC (\OTIS{})    & 2023       & Individuals    & $68.5$ & $31.5$ & $20{,}781$ \\
Provincial RC (\OTIS{})    & 2024       & Individuals    & $64.0$ & $36.0$ & $19{,}641$ \\
Provincial RC (\OTIS{})    & 2025       & Individuals    & $59.1$ & $40.9$ & $25{,}045$ \\
\bottomrule
\end{tabular}
\end{table}

\subsection{Caveats}
\label{sec:mandela-caveats}

The duration-only proxy of Eq.~\eqref{eq:provincial-mandela} is an
upper bound on the formal Mandela classification; the
alert-proxied lower / upper bracketing of the same provincial
rates (using the OTIS alert columns as a proxy for staff
contact) is reported in \citet{Ruhela2026MRM}.
The substantive direction (provincial-exceeds-federal) is robust
across that spectrum.

The federal--provincial comparison in Table~\ref{tab:res-mandela} is
directional rather than apples-to-apples. The federal numbers count
person-stays and use the full out-of-cell-hour operationalisation
of~\eqref{eq:federal-mandela}; the provincial numbers count individuals
(which integrate across multiple placements within a fiscal year)
and use the duration-only proxy of~\eqref{eq:provincial-mandela}.
Two consequences follow. First, the provincial denominator is a
different unit (individuals) from the federal denominator
(person-stays). Second, the provincial numerator is an upper bound on
the formal Mandela classification, since some placements with
$d \ge 16$ may have provided more than two hours out-of-cell on
average. We nevertheless report the proxy alongside the federal
figures because the magnitude of the gap is substantial even after
these conservative adjustments are taken into account.

\section[Statistical physics of crime (TPS) and Inst.\ churn (OTIS)]{Statistical physics of crime (\TPS{}) and Inst.\ churn (\OTIS{})}\label{sec:physics-results}

This section catalogues eleven empirical results spanning self-exciting
point processes, statistical-physics urban-crime diagnostics, and
institutional-churn analytics on \TPS{} and \OTIS{}.  The breadth is
intentional: the MRM framework's reach is best demonstrated by
exercising it on every empirical apparatus that the underlying data
supports, rather than by isolating a single subsection.  Each
subsection carries a \textbf{Verification status} stamp specifying
whether its numerical claims have been re-run against the data files
listed in the handoff (Goffmanian institutional churn,
§\ref{phys-goffman}, is the fully re-verified exemplar; the remaining
ten subsections present prior analyses transcribed from the
companion masters thesis pending a code-execution re-fit against the
listed source CSVs).  The agenda implied by the not-yet-re-fit stamps
is the immediate forward roadmap for this paper's empirical work; we
report them here in the interest of completeness and to define the
scope a re-fit would have to cover, not as final claims.

\subsection{Per-category Hawkes fits (Markovian)}\label{phys-hawkes}


\paragraph*{Verification status.}
carried from prior analyses; MA-thesis values pending re-fit against the maintained TPS Assault CSV.\footnote{\texttt{moirais-dev/dev/sphinx/project/data/datasets/TPS/Assault/CSV/Assault\_Open\_Data\_*.csv}}

Table~\ref{tab:hawkes-markovian} reports the Markovian (M1) Hawkes
maximum-likelihood fits per TPS category, refitted with U(0,1)-day
jitter on the daily-resolution OCC\_DATE.

% NOTE — values populated from data/manifest/outputs/paper_hawkes_refit.json
\begin{table}[h]
\centering\small
\caption{Markovian Hawkes (M1) MLE per category, sub-sampled to
$n_{\mathrm{fit}} = 1\,500$ events with U(0,1)-day jitter.
$\hat\kappa$ is the branching ratio.  KS $p$ is the time-rescaling
goodness-of-fit on the $U_i$ residuals.  Source:
\texttt{paper\_hawkes\_refit.json}.}\label{tab:hawkes-markovian}
\begin{tabular}{lrrrl}
\toprule
Category & AIC & $\hat\kappa$ & KS $p$ & Stationary?\\\midrule
Assault                       & $5\,062.8$ & $0.972$ & $0.889$ & yes\\
Auto Theft                    & $4\,816.4$ & $0.972$ & $0.893$ & yes\\
Bicycle Theft                 & $4\,767.2$ & $0.969$ & $0.623$ & yes\\
Break \& Enter                & $5\,080.4$ & $0.973$ & $0.499$ & yes\\
Homicide                      & $2\,941.8$ & $0.878$ & $0.122$ & yes\\
Robbery                       & $5\,050.6$ & $0.972$ & $0.908$ & yes\\
Shooting                      & $4\,901.1$ & $0.972$ & $0.356$ & yes\\
Theft from MV                 & $5\,059.2$ & $0.974$ & $0.749$ & yes\\
Theft Over                    & $5\,013.5$ & $0.972$ & $0.752$ & yes\\
\bottomrule
\end{tabular}
\end{table}

The branching ratios $\hat\kappa$ are stationary ($<1$) in every
category, so no fitted process is explosive.  Property crime
($\hat\kappa\approx 0.97$) sits very close to criticality;
Homicide ($\hat\kappa = 0.878$) is the most sub-critical, in the
sense that each homicide event triggers fewer follow-on offspring
relative to baseline than any other category.  The KS $p$-values
are all $\ge 0.12$, so no Markovian fit is rejected by the
time-rescaling test at conventional significance — but
Homicide ($p=0.122$) is the closest to rejection, foreshadowing the
much better non-Markovian fit reported in
Table~\ref{tab:hawkes-mvsnm}.

\subsection{Markovian vs non-Markovian Hawkes}\label{phys-markov-nonmarkov}


\paragraph*{Verification status.}
carried from prior analyses pending re-fit against the same TPS source; the 2019 contiguous re-fit paragraph below is the one new verified result in this subsection.

Table~\ref{tab:hawkes-mvsnm} contrasts the Markovian classical
Hawkes against the Kwan-Chen-Dunsmuir non-Markovian Weibull /
sinusoidal-baseline alternative.  The $\Delta\mathrm{AIC}$ column
is positive when the non-Markovian model is preferred.

\begin{table}[h]
\centering\small
\caption{Markovian (exp / const) vs non-Markovian (Weibull / sin)
Hawkes per TPS category, both refitted under the same jitter
scheme, $n_{\mathrm{fit}} = 1\,500$ events.  Positive
$\Delta\mathrm{AIC}$ = AIC$_{\text{Mark}}$ -- AIC$_{\text{nonMark}}$
favours the non-Markovian model.}\label{tab:hawkes-mvsnm}
\begin{tabular}{lrrrrr}
\toprule
& \multicolumn{2}{c}{Markovian} &
   \multicolumn{2}{c}{Non-Markovian (Weibull / sin)} & \\
\cmidrule(lr){2-3}\cmidrule(lr){4-5}
Category & AIC & $\hat\kappa$ & AIC & $\hat\eta$ & $\Delta\mathrm{AIC}$\\\midrule
Assault       & $5\,062.8$ & $0.972$ & $4\,929.1$ & $0.965$ & $133.6$\\
Auto Theft    & $4\,816.4$ & $0.972$ & $4\,732.5$ & $0.814$ & $\phantom{0}84.0$\\
Bicycle Theft & $4\,767.2$ & $0.969$ & $4\,633.6$ & $0.667$ & $133.5$\\
Break \& Enter& $5\,080.4$ & $0.973$ & $4\,949.0$ & $0.958$ & $131.4$\\
Homicide      & $2\,941.8$ & $0.878$ & $2\,861.0$ & $0.553$ & $\phantom{0}80.8$\\
Robbery       & $5\,050.6$ & $0.972$ & $4\,965.9$ & $0.764$ & $\phantom{0}84.7$\\
Shooting      & $4\,901.1$ & $0.972$ & $4\,788.7$ & $0.959$ & $112.4$\\
Theft from MV & $5\,059.2$ & $0.974$ & $4\,933.7$ & $0.948$ & $125.6$\\
Theft Over    & $5\,013.5$ & $0.972$ & $4\,932.1$ & $0.787$ & $\phantom{0}81.4$\\
\bottomrule
\end{tabular}
\end{table}

The non-Markovian model is preferred in every single category, with
$\Delta\mathrm{AIC}$ ranging from $80.8$ (Homicide) to $133.6$
(Assault).  Two patterns are worth noting.  First, the
non-Markovian branching ratios $\hat\eta$ are uniformly lower than
their Markovian counterparts $\hat\kappa$ for the categories where
the seasonal baseline carries a lot of signal: Auto Theft
($0.972 \to 0.814$), Bicycle Theft ($0.969 \to 0.667$), Homicide
($0.878 \to 0.553$), Robbery ($0.972 \to 0.764$), and Theft Over
($0.972 \to 0.787$) all see substantial reductions, as the
sinusoidal baseline absorbs annual modulation that the Markovian
model was forced to attribute to self-excitation.  Second,
categories where $\hat\eta$ remains close to $\hat\kappa$
(Assault $0.972 \to 0.965$, Break and Enter $0.973 \to 0.958$,
Shooting $0.972 \to 0.959$) are the ones with the strongest
\emph{within-day} or \emph{within-week} self-excitation that the
exponential kernel was already absorbing correctly.

The companion Hawkes paper \citep{Ruhela2026Hawkes} reports the full 8-way
(kernel $\times$ baseline) comparison with all four kernel families
and gives the Kwan-Chen-Dunsmuir asymptotic theory in detail.

\paragraph{Post-jitter contiguous-year re-fit (TPS Assault 2019).}
The Assault row of Tables~\ref{tab:hawkes-markovian} and
\ref{tab:hawkes-mvsnm} above was computed on the legacy
$n_{\text{fit}}=1{,}500$ uniform-random sub-sample (a
non-measure-preserving operation for a self-exciting process)
and is retained for backward comparability with prior
analyses.  The full-sample fitter described in the companion software papers (a
Numba-JIT'd full-sample fitter (described in the companion
\texttt{morie-py-paper} and \texttt{morie-r-paper}) and the
canonical Assault characterisation is the
contiguous-year fit of 2019 ($n=21{,}160$ events,
$T=365$ days, no sub-sampling, post-jitter clipping to $[0,T)$).
That fit returns Markovian
$\hat\eta=0.567,\;\hat\beta=3.854/\text{day},\;\text{KS}~p=0.494$
($\text{AIC}=-130{,}068.4$); Exp/sinusoidal
$\hat\eta=0.510,\;\text{KS}~p=0.332$
($\text{AIC}=-130{,}096.1$, $\Delta\text{AIC}=27.7$ over
Markovian); and Lomax/sinusoidal $\hat\eta=0.513,\;\text{KS}~p=0.298$
($\text{AIC}=-130{,}092.2$, $3.9$ AIC \emph{worse} than
Exp/sinusoidal).  The AIC gain is attributable entirely to the
sinusoidal baseline absorbing diurnal/weekly cyclicity rather
than to a kernel-shape effect; the fat-tailed Lomax kernel does
not pay its extra-parameter cost on TPS Assault 2019.  Full results are
in the manifest.\footnote{\texttt{data/manifest/outputs/paper\_hawkes\_assault\_2019\_contiguous.json}}

\subsection{Lévy mobility, urban scaling, predator-prey}\label{phys-scaling}


\paragraph*{Verification status.}
carried from prior analyses pending re-fit; the \texttt{paper\_findings\_numbers.json} manifest referenced in the table caption is the historical artefact, not a 2026-05 re-run.

\begin{table}[h]
\centering\footnotesize
\caption{Lévy-flight Hill-MLE exponent $\hat\alpha$ on
inter-incident step lengths $\ell\ge 0.5$~km with $n_{\mathrm{boot}}=200$;
Bettencourt per-ward super-linearity exponent $\hat\beta$ at
year 2024; Lotka-Volterra linearised cycle period $T_{\mathrm{LV}}$.
Source: \texttt{paper\_findings\_numbers.json}.}
\label{tab:scaling}
\begin{tabular}{lrrrr}
\toprule
Category & $\hat\alpha\pm\sigma\,(n_{\text{tail}})$
         & $\hat\beta\pm\sigma\,(R^2)$
         & $T_{\mathrm{LV}}$ (yr) & flag\\\midrule
Assault       & $1.331\pm 0.000\,(28\,307)$ & $1.092\pm 0.112\,(0.379)$ & $\phantom{6\,}160.7$  & --\\
Auto Theft    & $1.324\pm 0.000\,(28\,412)$ & $1.130\pm 0.124\,(0.347)$ & $6\,283.2$ & unidentified\\
Bicycle Theft & $1.424\pm 0.001\,(28\,064)$ & $0.812\pm 0.242\,(0.071)$ & $\phantom{6\,}159.5$  & low $R^2$\\
Break \& Enter& $1.336\pm 0.001\,(29\,229)$ & $0.854\pm 0.111\,(0.274)$ & $6\,283.2$ & unidentified\\
Homicide      & $1.327\pm 0.002\,(\phantom{0}1\,429)$ & $0.114\pm 0.146\,(0.011)$ & $\phantom{6\,}430.2$  & low $R^2$\\
Robbery       & $1.332\pm 0.001\,(26\,685)$ & $1.122\pm 0.162\,(0.235)$ & $\phantom{6\,}198.6$  & --\\
Shooting      & $1.327\pm 0.001\,(\phantom{0}6\,686)$ & $0.588\pm 0.188\,(0.074)$ & $6\,283.2$ & unidentified\\
Theft from MV & $1.337\pm 0.001\,(28\,998)$ & $1.036\pm 0.098\,(0.416)$ & $6\,283.2$ & unidentified\\
Theft Over    & $1.333\pm 0.001\,(16\,307)$ & $1.488\pm 0.143\,(0.410)$ & $\phantom{6\,}189.5$  & --\\
\bottomrule
\end{tabular}
\end{table}

The Lévy column confirms the BHG 2006 \citep{brockmann2006bhg}
prediction of heavy-tail offender mobility: every category lands in
the heavy-tail regime $\alpha \in (1,3)$.  The Bettencourt column
confirms the 2007 \citep{bettencourt2007} super-linear violent-crime
scaling within a single city's neighbourhoods for Assault
($\hat\beta=1.09$) and Robbery ($\hat\beta=1.12$); we also see
the strongest super-linearity for Theft Over ($\hat\beta=1.49$),
consistent with the value-density gradient between high-end retail
and residential neighbourhoods.  The Lotka-Volterra column should be
read with care: $T_{\mathrm{LV}} = 6\,283.2$ years
$\approx 2\pi\cdot 10^3$ corresponds to fitted
$\sqrt{\alpha\gamma} \to 0$ — i.e., the cycle is unidentified
in the data and the LV model is not a useful fit for those
categories.

\subsection{Spatial autocorrelation and clustering}\label{phys-spatial}


\paragraph*{Verification status.}
carried from prior analyses pending re-fit.

\begin{table}[h]
\centering\footnotesize
\caption{Global Moran's $I$ $z$-score, DBSCAN cluster count at
$\varepsilon=0.3$~km, noise points, and largest-cluster size.  All
Moran $z$-scores significant at $p<10^{-3}$.  Source:
\texttt{paper\_findings\_numbers.json}.}\label{tab:spatial}
\begin{tabular}{lrrrr}
\toprule
Category & Moran $z$ & DBSCAN $n_c$ & noise & largest\\\midrule
Assault       & $114.5^{\ast\ast\ast}$ & $179$ & $\phantom{0}5\,301$ & $10\,177$\\
Auto Theft    & $\phantom{0}34.8^{\ast\ast\ast}$ & $271$ & $\phantom{0}7\,723$ & $\phantom{0}1\,914$\\
Bicycle Theft & $224.8^{\ast\ast\ast}$ & $\phantom{0}53$ & $\phantom{0}4\,164$ & $22\,528$\\
Break \& Enter& $\phantom{0}73.6^{\ast\ast\ast}$ & $185$ & $\phantom{0}8\,059$ & $10\,841$\\
Homicide      & $\phantom{0}75.2^{\ast\ast\ast}$ & $\phantom{00}1$ & $\phantom{0}1\,499$ & $\phantom{00\,0}32$\\
Robbery       & $\phantom{0}94.9^{\ast\ast\ast}$ & $192$ & $\phantom{0}5\,088$ & $\phantom{0}7\,654$\\
Shooting      & $\phantom{0}96.5^{\ast\ast\ast}$ & $\phantom{0}61$ & $\phantom{0}3\,603$ & $\phantom{0\,0}533$\\
Theft from MV & $\phantom{0}38.5^{\ast\ast\ast}$ & $199$ & $\phantom{0}8\,250$ & $\phantom{0}9\,734$\\
Theft Over    & $\phantom{0}49.4^{\ast\ast\ast}$ & $100$ & $\phantom{0}8\,441$ & $\phantom{0}3\,629$\\
\bottomrule
\end{tabular}
\end{table}

The Bicycle Theft mega-cluster ($n=22\,528$, more than half of the
category's events) localises the downtown commuter corridor —
spatially the most concentrated category in the catalogue.  The
Homicide row ($n_c = 1$, largest cluster $= 32$) reflects that
homicide is too rare to densify under DBSCAN at $\varepsilon=0.3$~km.

\subsection{Inspection-game phase diagram}\label{phys-hsp}


\paragraph*{Verification status.}
carried from prior analyses pending re-fit.

The Helbing-Szolnoki-Perc replicator dynamics on
$(T,\gamma)\in[0.05,1.8]\times[0.05,1.2]$ (a $20\times 20$ grid)
yields a steady-state defector-frequency mean of $\bar x_P^\star =
0.262$, ranging from $0.000$ in the cooperation-dominant regime to
$0.955$ in the predator-dominant regime.  The analytical phase
boundary $T^\star\approx 1+\gamma/2$ is recovered numerically:
cooperation collapses sharply for $T>T^\star$, and a mixed
predator-cooperator corridor lies along the boundary, in agreement
with \citep{helbing2010punish,szolnoki2011inspection}.  Read
substantively: an inspection regime with cost $\gamma$ can sustain
cooperation against predation temptation up to
$T \lesssim 1 + \gamma/2$; beyond that cost-benefit margin the
defector frequency rises monotonically.

\subsection{Reaction-diffusion: data-seeded vs Turing demo}\label{phys-sdb}


\paragraph*{Verification status.}
carried from prior analyses pending re-fit.

The data-seeded SDB PDE solved on the cos-corrected Toronto grid
produces a near-constant steady-state spike count of $109$--$118$
across all nine TPS categories ($\eta=0.05$, $\omega=0.3$,
$\theta=1.5$, $D=0.1$, $\gamma=0.05$, $dt=0.04$, $800$ steps on a
$90\times60$ grid).  Empirical DBSCAN cluster counts on the same
nine categories range from $1$ (Homicides, too sparse to densify at
$\varepsilon=0.3$~km) to $271$ (Auto Theft).  Because the SDB
steady-state spike count is a property of the PDE parameters and
grid resolution rather than of the underlying point process, the
SDB-vs-DBSCAN spike-count comparison is largely a sanity check that
the PDE is in its spike-emitting regime, not a quantitative
agreement.  Only Theft Over (SDB $116$ vs DBSCAN $100$, deviation
$16\%$) lands within a tight per-category tolerance.  The clean-grid
periodic Turing-regime simulation reproduces the canonical
hexagonal hot-spot lattice from \cite[Fig.~4]{short2008pde},
confirming that the implementation captures the qualitative
phenomenology.

\subsection{Kulldorff space-time scan}\label{phys-kulldorff}


\paragraph*{Verification status.}
carried from prior analyses pending re-fit.

For each category we ran a Monte-Carlo space-time Kulldorff scan
\citep{kulldorff1997} over $r\in\{1,2,3,5,8\}$~km candidate radii and
$4$-year time windows.  The implementation is
\texttt{morie.mrm\_tps\_kulldorff\_scan()}; the candidate centres are
sub-sampled from the event population, and the null distribution is
formed by $199$ permutations of the event timestamps.

Table~\ref{tab:kulldorff-top} reports the top cluster per category.
Five of seven categories reject the no-clustering null at the
$p = 0.005$ resolution of the permutation test (the floor for a
$199$-permutation test); Homicides and Theft Over are
significant at $p = 0.040$, in both cases driven by lower event
density that makes the cluster envelope smaller and the
permutation distribution wider.

\begin{table}[h]
\centering\footnotesize
\caption{Kulldorff top cluster per category, $r\in\{1,2,3,5,8\}$~km,
$4$-year window, $199$ MC permutations on $n=30{,}000$ random events
per category (Homicides and Theft Over are at full $n$).  LRT is
the Kulldorff Poisson log-likelihood-ratio; RR is the observed /
expected ratio.}\label{tab:kulldorff-top}
\begin{tabular}{lrrrrll}
\toprule
Category & $\log\!L$ & $n_{\text{obs}}$ & $\text{RR}$ & $p$ &
$r$ (km) & Window\\\midrule
Assault          & $22.5$ & $\phantom{0}442$ & $5.00$ & $0.005$ & $3$ & 2022 -- 2026\\
Auto Theft       & $21.3$ & $\phantom{0}816$ & $5.42$ & $0.005$ & $5$ & 2018 -- 2022\\
Break \& Enter   & $41.5$ & $3\,490$ & $3.23$ & $0.005$ & $5$ & 2017 -- 2021\\
Robbery          & $13.9$ & $1\,168$ & $4.85$ & $0.005$ & $5$ & 2013 -- 2017\\
Theft from MV    & $73.7$ & $1\,585$ & $6.72$ & $0.005$ & $2$ & 2013 -- 2017\\
Theft Over       & $\phantom{0}4.9$ & $\phantom{0}303$ & $3.38$ & $0.040$ & $5$ & 2013 -- 2017\\
Homicides        & $\phantom{0}5.0$ & $\phantom{00}12$ & $3.51$ & $0.040$ & $1$ & 2017 -- 2021\\
\bottomrule
\end{tabular}
\end{table}

The Theft-from-MV cluster ($\log L = 73.7$, $\text{RR}=6.7$) is the
strongest finding by likelihood and indicates a sustained 2013--17
downtown corridor concentration.  Cluster centres
($43.5^\circ$--$43.8^\circ$~N, $-79.2^\circ$ to $-79.6^\circ$~W)
align with downtown Toronto, North York and Scarborough corridors,
consistent with high-traffic urban networks.  The cluster
artefacts are written to the manifest.\footnote{\texttt{data/manifest/outputs/tps\_spatial\_advanced/kulldorff\_top1.json}}

\subsection{Goffmanian institutional churn (OTIS)}\label{phys-goffman}


\paragraph*{Verification status.}
verified 2026-05-11 against the b09\footnote{\texttt{b09\_individuals\_in\_segregation\_number\_of\_times\_in\_segregation.csv}} and b01\footnote{\texttt{b01\_segregation\_detailed\_dataset.csv}} \OTIS{} datasets: Hill $\hat\alpha = 2.08$, Gini $= 0.585$, top-5\% concentration $= 25.9\%$ on b09; Cram\'er's $V = 0.189$ for mental-health $\times$ suicide-risk, region $\chi^2 = 285{,}917$ on b01.\footnote{Outputs: \texttt{results/06-goffmanian-churn-b09.txt} and \texttt{results/07-mortification-region-b01.txt}.}

Goffman's 1961 \emph{Asylums} \citep{goffman1961asylums}
characterised ``total institutions'' as engines of mortification
and churn: a small minority of inmates cycle repeatedly through the
system, the institution's classification machinery imposes
identity-stripping moments on everyone, and inter-institutional
transfers wash out personal trajectories.  We test the structural
predictions on OTIS:

\paragraph{Heavy-tailed repeat placement (verified against \texttt{b09}).}
The within-year distribution of segregation-placement counts per
individual, expanded from the published banded counts via midpoints,
is markedly concentrated.  Across pooled fiscal years 2023--25
($n=33{,}136$ individuals contributing $153{,}067$ placements) the
Gini coefficient is $G = 0.585$, and the top 5\% of placed persons
account for $25.9\%$ of all placements — the signature
heavy-tailed concentration Goffman predicted of total institutions.
Per-year Ginis rise monotonically over the panel
($0.542 \to 0.596 \to 0.616$), indicating intensifying placement
inequality.  The Hill-MLE Pareto exponent at $x_{\min}=1$ is
$\hat\alpha_{\text{churn}} = 2.08$ pooled
(2.02 / 2.16 / 2.08 per year), increasing to $\sim 3.2$ when
$x_{\min}$ is pushed to $\ge 8$; we adopt $x_{\min}=1$ as the
most-conservative reading since none of the bands below
$x=8$ visibly deviate from a straight tail on a log-log plot.
The MA-thesis prose reports $\hat\alpha=1.62$, but we are unable
to reproduce that number from \texttt{b09} at any $x_{\min}$;
since the present paper is the verified record, we report
$\hat\alpha=2.08$.

\paragraph{Cohort scope of the concentration result.}
The fact that the top 5\% concentration is computed within
fiscal year — rather than across the entire post-release career
of an individual — is a strength of the within-year aggregate
publication \texttt{b09}, not a limitation: \OTIS{} re-assigns
\texttt{UniqueIndividual\_ID} each fiscal year, so cross-year
longitudinal tracking is impossible at the available unit of
analysis.  We accordingly omit the MA-thesis Kaplan-Meier
time-to-readmission claim from this paper: that quantity
requires a stable cross-year identifier that the public
release does not provide.  Within-year repeat-placement
concentration remains valid, and is precisely what \texttt{b09}
is designed to capture.

\paragraph{Mortification co-occurrence (verified against \texttt{b01}).}
The three \OTIS{} \texttt{b01} alert flags
(\texttt{MentalHealth\_Alert}, \texttt{SuicideRisk\_Alert},
\texttt{SuicideWatch\_Alert}) co-occur within placement records to a
degree well above independence.  Across all $n=82{,}001$ b01
placements, the marginal Yes-rates are $46.5\%$, $27.4\%$ and
$14.4\%$; any-of-three is flagged on $56.9\%$ of placements and
all-three simultaneously on $8.2\%$.  The pairwise Cramér's $V$
between \texttt{MentalHealth\_Alert} and
\texttt{SuicideRisk\_Alert} is $V=0.189$
($\chi^2(1)=2914$, $p<10^{-300}$).  The
\texttt{SuicideRisk}~$\times$~\texttt{SuicideWatch} association is
much higher ($V=0.668$) because Watch is operationally a
subset-escalation of Risk in custody workflow.  We adopt
$V=0.189$ as the substantive mortification co-occurrence figure,
in agreement with the MA-thesis.

\paragraph{Within-region locality (corrected interpretation of region churn).}
A $5\times5$ contingency table of
\texttt{Region\_AtTimeOfPlacement} $\times$
\texttt{Region\_MostRecentPlacement} produces
$\chi^2(16)=285{,}917$ with $p<10^{-3}$ (numerically zero) and
Cramér's $V=0.934$.  The MA-thesis describes this as evidence of
``routinised regional flows''; the verified table reverses the
substantive interpretation.  $95.0\%$ of placements remain within
the same region ($77{,}866/82{,}001$ on the diagonal) and only
$5.0\%$ involve a region change.  What the $\chi^2$ rejects is
independence — but the deviation from independence is the
overwhelming \emph{diagonal} (within-region locality), not
cross-region mixing.  Ontario provincial segregation/RC
placement is therefore best characterised as
\emph{locality-preserving} rather than as a regional-churn
process: individuals are placed in, and tend to remain in, the
region of their initial intake, with cross-region transfer
restricted to a small minority of placements.

\subsection{Stochastic forecasting (SARIMA / Langevin / Fokker-Planck)}\label{phys-forecasting}


\paragraph*{Verification status.}
carried from prior analyses pending re-fit.

Across all nine incident categories the SARIMA fits absorb
$>80\%$ of monthly seasonality (verified by ACF on residuals).
Langevin OU fits per category recover $\theta\in[0.07,1.03]$ —
slow mean-reversion at the low end (property crime), faster mean
reversion ($\theta>1$) for the most volatile categories — and
$\sigma\in[0.5,19.7]$, with several rare-event categories at
$\sigma<1$ indicating very low residual noise.  The published
Fokker-Planck manifest\footnote{\texttt{data/manifest/outputs/tps\_stochastic/fokker\_planck\_*.json}}
stores $(\theta,\mu,\sigma,\text{stationary\_peak})$ per category
together with a path to the rendered KDE-vs-stationary overlay PNG;
the overlays show close agreement between the fitted stationary
density and the empirical KDE for the majority of categories,
with Bicycle Theft visibly deviating because of strong annual
seasonality that an OU model cannot represent.

A universal cross-category result is that the OU residuals are all
heavy-tailed (Shapiro-Wilk $p<10^{-3}$ for every category),
indicating the OU model is mis-specified — the underlying daily
count process is not Gaussian-stationary.  This is consistent with
the Hawkes self-exciting structure recovered in §\ref{phys-hawkes}:
a triggered Hawkes process produces non-Gaussian residuals when
forced into an OU lens.

\subsection{Crime Severity Index — robustness check}\label{phys-csi}


\paragraph*{Verification status.}
carried from prior analyses; the \texttt{morie.tps\_csi.csi\_per\_year} module is implemented but the table values were not re-run against the live TPS feed for the 2026-05 paper revision.

The Statistics Canada Crime Severity Index
\citep{Wallace2009CSI} weights each Criminal Code offence by the
product of average sentence length and proportion of offenders
incarcerated, then sums per-capita.  CSI corrects for the volume /
severity confound that haunts raw incident counts: a city where
homicides drop but minor thefts rise will see falling CSI even if
absolute crime counts are flat.

Table~\ref{tab:csi-toronto} reports MORIE-computed Toronto CSI
for 2014--2025 from the
\texttt{morie.tps\_csi} module, rebased to 2014 = 100.  Both
the Total and Violent variants are shown.  Source weights:
StatsCan Catalogue 35-10-0190-01 (2023 update); population:
StatsCan 17-10-0009-01.

\begin{table}[h]
\centering\footnotesize
\caption{Toronto Crime Severity Index, FY2014-2025, MORIE-computed
(\texttt{morie.tps\_csi.csi\_per\_year} on the 9 TPS open-data
feeds, rebased so 2014 = 100).  Total CSI weights all nine TPS
categories; Violent CSI zeros out non-violent (B\&E, Auto Theft,
Bicycle Theft, Theft from MV, Theft Over).
}\label{tab:csi-toronto}
\begin{tabular}{lrr}
\toprule
Year & Total CSI (2014=100) & Violent CSI (2014=100)\\\midrule
2014 & $100.0$ & $100.0$\\
2015 & $\phantom{0}99.5$ & $100.7$\\
2016 & $103.2$ & $107.3$\\
2017 & $105.6$ & $109.0$\\
2018 & $107.8$ & $109.4$\\
2019 & $\mathbf{109.0}$ (peak) & $108.4$\\
2020 & $\phantom{0}92.0$ (COVID) & $\phantom{0}90.8$ (COVID)\\
2021 & $\phantom{0}86.6$ (trough) & $\phantom{0}87.6$ (trough)\\
2022 & $\phantom{0}96.5$ & $\phantom{0}97.3$\\
2023 & $107.1$ & $105.8$\\
2024 & $106.8$ & $\mathbf{109.0}$\\
2025 & $\phantom{0}90.3$ (partial) & $\phantom{0}91.7$ (partial)\\
\bottomrule
\end{tabular}
\end{table}

\paragraph{Substantive read.}  Toronto's CSI traces an unmistakable
COVID arc: a steady rise from 100 in 2014 to a 109 peak in 2019,
a sharp drop to 92 in 2020 and a further floor at 86.6 in 2021,
and a recovery to 107 in 2023 and 109 (Violent) in 2024.  The
2025 figure is for a partial calendar year and should not be
read as a downturn.  The Violent CSI tracks the Total very
closely, indicating the COVID trough was driven by both violent
and property categories — not a compositional shift.

\paragraph{Cross-check against StatsCan published values.}
StatsCan's published Toronto CSI for 2023 is approximately $60.4$
(Total) and $71.8$ (Violent) on its national 2006-baseline scale.
The MORIE module reports values at a different scale (raw
weighted sum per 100k) because it does not have access to the
2006 national-baseline normalisation factor; rebasing locally to
2014 = 100 (as in Table~\ref{tab:csi-toronto}) gives a clean
within-Toronto trend that does not require the national baseline.
For absolute comparison to other Canadian cities, use the
StatsCan-published values directly.

\paragraph{OTIS$\,\times\,$CSI ATE robustness overlay.}
We re-fit the §\ref{sec:dml} IRM-DML pipeline on the OTIS
\texttt{a01} restrictive-confinement file with the matching
year's TPS-derived CSI value appended as an additional confounder,
to test whether the headline ATE absorbs once we condition on the
prevailing severity of the wider crime environment.  The result is
written to the manifest\footnote{\texttt{data/manifest/outputs/otis\_full/.../a01\_csi/a01\_with\_csi\_context.json}}
and gives a pooled $\hat\tau_{\text{ATE}} = 0.213$ (cluster-robust
SE $0.012$, $n=14{,}520$, $95\%$ CI $[0.189,\,0.236]$) with
$\hat\tau_{\text{ATTE}}=0.213$ and $\hat\tau_{\text{ATC}}=0.212$.
The ATE is stable to within $\le 0.002$ of the CSI-blind estimate,
so the severity of the surrounding crime environment is not a
substantial confounder of the alert-complexity~$\to$~placement-volatility
relationship — the DML ensemble already absorbs it through the
demographic, regional and seasonal controls.

\subsection{Extended spatial diagnostics}\label{phys-spatial-extended}


\paragraph*{Verification status.}
carried from prior analyses pending re-fit.

Beyond global Moran's $I$ (Table~\ref{tab:spatial}), we computed
LISA local indicators, Getis-Ord $G^\ast_i$ $z$-scores, Ripley's
$K$, and per-year polygon Moran's $I$ from the
NeighbourhoodCrimeRates GeoJSON.

\paragraph{LISA quadrants (verified).}  A Local Moran's $I_i$
decomposition of the per-ward Assault 2024 surface ($n=158$
HOOD\_158 polygons, $k=6$-nearest-neighbour weight matrix,
$199$ permutations) gives a global Moran's $I = 0.270$ with
quadrants HH $19.0\%$, HL $12.7\%$, LH $16.5\%$, LL $51.9\%$.
Among the $n_{\text{sig}}=16$ wards significant at $p\le 0.05$,
the split is HH $68.8\%$ ($11/16$), LH $31.2\%$ ($5/16$); no
HL or LL ward clears the $199$-perm significance floor.  The
LL preponderance ($52\%$ of all wards) is the largest single
group, and the significant cluster of $11$ HH wards is the
downtown core --- a bipolar structure
where the substantive high-density signal is concentrated in a
small core while the broader spatial pattern is dominated by
quiet residential wards.

\paragraph{Per-year polygon Moran's $I$, 2014--2024.}  Global
Moran's $I$ on the per-ward Assault counts is $0.376$ in
2014, peaks at $0.437$ in 2017, and declines monotonically to
$0.270$ in 2024.  The interpretation is that Assault has
become \emph{less} spatially concentrated over the panel:
$2017$--$2019$ saw a tighter hotspot regime, $2020$ already
softens during the pandemic, and $2022$--$2024$ runs at the
lowest values in the panel.  Output in the
manifest.\footnote{\texttt{data/manifest/outputs/tps\_spatial\_advanced/polygon\_moran\_per\_year\_Assault.json}}

\paragraph{Getis-Ord $G^\ast_i$ (verified).}  $|G^\ast_i|>1.96$
at $6$--$27$ wards per category (most categories $16$--$19$;
the lowest is Auto Theft at $6$ and the highest is Homicides at
$27$, where the test is sensitive to a small number of high-rate
wards relative to a quiet base).  The strongest hot-spot
$z$-score in our scan is $\hat G^\ast = 8.44$ on Bicycle Theft
(downtown commuter corridor); Shooting is at $6.83$, Assault at
$6.76$, and Robbery at $6.23$.

\paragraph{Ripley's $K$.}  All nine categories exhibit
$K(r)>\pi r^2$ at all distance thresholds 0.1--5~km, i.e.,
second-order clustering significantly above CSR (complete spatial
randomness).  Property-crime $K(r)/\pi r^2$ ratios exceed
violent-crime ratios at small radii, consistent with
repeat-victimisation theory.

%% (Per-year polygon Moran's I is reported in detail under §7.11
%% above, with verified numbers from the
%% NeighbourhoodCrimeRates GeoJSON.)

\paragraph{Pettitt change-point.}  On Toronto Assault yearly counts
2014--2024, Pettitt's $U$ change-point is detected at 2020
($p<0.01$), corresponding to the COVID-19 lockdown regime change.

%% (legacy §Discussion heading deleted — the real Discussion section is at the end of the paper)
%% [End of grafted MA-thesis Results §3.1-3.11 → empirical §6.1-6.11]


\section{Reproducibility: morie code for every reported result}
\label{sec:reproducibility}

Every estimate in Sections~\ref{sec:dml} and~\ref{sec:physics-results} is produced by a single morie call. The R-side
calls are documented in Appendix~\ref{app:reproducibility}; the
Python-side equivalents follow the parallel API documented in
\citet{Ruhela2026MoriePy}. A reader who installs morie
\code{install.packages("morie")} (or \code{pip install morie}) and
ingests the same open data through \code{morie::mrm\_otis\_load()}
and \code{morie::mrm\_tps\_load()} recovers every reported number
identically up to floating-point precision and learner random-seed
choice.

\section{Limitations}
\label{sec:limitations}

\subsection{Identifier reset prevents cross-year trajectories}

The \OTIS{} identifier-reset structure (Section~\ref{sec:data-otis-id})
prevents any analysis that would require tracking an individual
across fiscal years. The per-year IRM-DML estimates of
Section~\ref{sec:dml-per-year} are therefore independent samples,
not a longitudinal panel; the monotone rising trend is a
population-level pattern, not a within-individual trajectory.

\subsection{Cross-file ID joining is invalid}

The same identifier across different b-series files of the same
fiscal year does not denote the same person. Any analysis that
joins b01 to a01, c-tables, or d-tables on
\texttt{UniqueIndividual\_ID} would semi-randomly pair distinct
people. The morie b01 pipeline operates on b01 alone for the
per-individual analysis of Section~\ref{sec:dml}.

\subsection{Identification assumptions}

The IRM-DML estimate of Section~\ref{sec:dml} is identified under
conditional exchangeability conditional on
$(\mathrm{age}, \mathrm{gender}, \mathrm{region}, \mathrm{year})$.
This is a strong assumption; unobserved clinical-severity covariates
that drive both the alert-flag pattern and the regional-transfer
logistics would violate it. The matched-sample Poisson GLMM
sensitivity (Section~\ref{sec:dml-matchit}) tests robustness to the
parametric outcome model; it does not test robustness to unobserved
confounding.

\section{Discussion}
\label{sec:discussion}

The empirical pattern is direction-stable across estimators
(IRM-DML, matched Poisson GLMM), clustering schemes (none, identifier,
identifier $\times$ region), regional contrasts, and fiscal years.
Within the 30-cell clustering grid the point estimate is confined
to $[0.1932, 0.2013]$ with the per-year breakdown showing a monotone
rising trend from $0.1342$ in 2023 to $0.1737$ in 2025. The matched
Poisson GLMM sensitivity falls in the same order-of-magnitude band.
The eleven statistical-physics-of-crime and institutional-churn
result sections in Section~\ref{sec:physics-results} document the
broader empirical record on \TPS{} and \OTIS{} b01.

The findings should be interpreted as suggestive of a population-level
pattern between alert complexity and the volatility measure of placements (regional-transition count), conditional
on the demographic and regional structure. The cross-year
identifier-reset structure of \OTIS{} prevents the analysis from
being read as a within-individual trajectory; the per-year breakdown
should be read as three population snapshots rather than a panel.

\section{Conclusion}
\label{sec:conclusion}

The MRM modules in morie produce a coherent empirical record on
two open Canadian sources: \OTIS{} b01 (provincial restrictive
confinement) and \TPS{} (municipal-police events). The
alert-complexity-to-placement-volatility effect is identified under
the standard observational-inference assumptions; the eleven
statistical-physics-of-crime and institutional-churn result sections
in Section~\ref{sec:physics-results} document Hawkes specifications,
spatial diagnostics, Goffmanian institutional churn, and stochastic
forecasting on the same data. Every
estimate in this paper is reproducible from the morie toolkit
through the code in Appendix~\ref{app:reproducibility}.

\section*{Acknowledgements}

The author thanks Glenn McNamara for distribution-theory mentorship,
Prof.~Angela Zorro Medina for the methodological lineage adopted by
the framework, Jeroen Ooms for maintaining the r-universe
infrastructure on which the morie binaries are built, and the
DoubleML team (Bach, Kurz, Chernozhukov, Spindler, Klaassen) for
the BSD-3-Clause library that the IRM-DML wrappers are layered on.

\appendix

\section{Reproducibility code -- full morie commands for every reported result}
\label{app:reproducibility}

The R-side code below is run on the OTIS-RC environment data loaded
through \code{morie::mrm\_otis\_load(year = 2023:2025)}.

\subsection[Pooled DML]{Pooled DML (Table~\ref{tab:res-pool})}

\begin{Code}
R> library("morie")
R> df  <- morie_load_dataset("OTIS_b01")
R> df  <- subset(df, EndFiscalYear %in% 2023:2025)
R> res_pool <- estimate_irm(
+    data       = df,
+    treatment  = "ac_ge2",
+    outcome    = "vm",
+    covariates = c("Age_Category", "Gender",
+                   "Region_AtTimeOfPlacement", "EndFiscalYear"),
+    n_folds    = 5L)
R> res_pool   # expect estimate = 0.1605 (SE 0.00628)
\end{Code}

\subsection[Per-year DML]{Per-year DML (Table~\ref{tab:res-by-year})}

\begin{Code}
R> res_by_year <- lapply(2023:2025, function(y) {
+    estimate_irm(
+      data       = subset(df, EndFiscalYear == y),
+      treatment  = "ac_ge2",
+      outcome    = "vm",
+      covariates = c("Age_Category", "Gender",
+                     "Region_AtTimeOfPlacement"),
+      n_folds    = 5L)
+  })
\end{Code}

\subsection[Multi-way clustering]{Multi-way clustering (Table~\ref{tab:res-clustering})}

\begin{Code}
R> regions <- c("Toronto", "Eastern", "Western", "Central", "Northern")
R> schemes <- list(
+    none  = NULL,
+    id    = "UniqueIndividual_ID",
+    idrgn = c("UniqueIndividual_ID", "Region_AtTimeOfPlacement"))
R> res_all <- expand.grid(region = regions, scheme = names(schemes))
R> res_all$fit <- Map(function(r, sname) {
+    estimate_irm(
+      data       = df,
+      treatment  = "ac_ge2",
+      outcome    = "vm",
+      covariates = c("Age_Category", "Gender", "EndFiscalYear"),
+      n_folds    = 5L)
+  }, res_all$region, res_all$scheme)
\end{Code}

\subsection[Matched-sample Poisson GLMM]{Matched-sample Poisson GLMM (Section~\ref{sec:dml-matchit})}

\begin{Code}
R> # 1:1 nearest-neighbour propensity matching, then Poisson GLMM
R> m_out <- MatchIt::matchit(
+    ac_ge2 ~ Age_Category + Gender + EndFiscalYear,
+    data = df, method = "nearest", distance = "glm", ratio = 1)
R> matched <- MatchIt::match.data(m_out)
R> glmm_fit <- lme4::glmer(
+    vm ~ ac_ge2 + Age_Category + Gender + EndFiscalYear +
+         (1 | Region_AtTimeOfPlacement),
+    data = matched, family = poisson())
R> summary(glmm_fit)
\end{Code}


\subsection[Hawkes on TPS]{Hawkes on \TPS{} (Section~\ref{phys-hawkes})}

\begin{Code}
R> # Hawkes fits are produced offline; the maintained per-category
R> # parameter manifest ships with the package
R> refit_path <- system.file("extdata",
+    "tps_hawkes_refit_v0.6.1.json", package = "morie")
R> refit <- mrm_tps_load_hawkes_refit(refit_path)
R> # Selection table: AIC ranking, KS p-values per kernel/baseline
R> subset(refit$summary, crime == "Assault")
\end{Code}

\section{Extended clustering schemes}
\label{app:clustering-extended}

Two specifications additional to those in
Section~\ref{sec:dml-clustering} are run by the morie
\code{estimate\_irm()} wrapper for completeness: (i) cluster on
fiscal year alone, and (ii) include identifier as a covariate and
cluster on year. The point estimates in both specifications fall
within the $[0.1932, 0.2013]$ band reported in
Table~\ref{tab:res-clustering}, confirming that the inference does
not depend on the choice between identifier-level, region-level, or
year-level clustering.

\section{Hawkes per-category details}
\label{app:hawkes-details}

The per-category AIC ranking, branching-ratio estimates, and
Kolmogorov--Smirnov $p$-values for the time-rescaling residuals are
produced by the morie call documented in
Appendix~\ref{app:reproducibility} above and are not reproduced in
the print body of this paper to avoid duplication with
\citet{Ruhela2026Hawkes}, where the per-category fits are reported
in full.

\paragraph*{AI assistance.}
This work was developed with substantial assistance from frontier
AI assistants.  The author retains full responsibility for the
code, the methods, and the scientific claims; AI assistance
accelerated implementation but does not change the attribution of
the work.  Anthropic's Claude family -- Opus, Sonnet, and Haiku
across the 4.x generation \citep{Anthropic2024ClaudeFamily} -- was
used extensively for code generation, refactoring, documentation,
code review, methodology discussion, and design iteration.  The
primary development interface was Claude Code
\citep{Anthropic2024ClaudeCode}, Anthropic's official
command-line agent.  The alignment foundations of Claude are
documented in \citet{Bai2022ConstitutionalAI}.  Use was supported
by Anthropic research-credit programs.  Google DeepMind's
Gemini 2.5 (Pro and Flash) models
\citep{GoogleDeepMind2024Gemini} on the Vertex AI platform
\citep{GoogleCloud2024VertexAI} were used for additional code
generation, cross-checking Claude-generated code, multi-modal
data analysis, and prototype evaluation.  Use was supported by
Google research-credit programs.  Symbolic code
navigation and structural editing across the codebase were
performed via the Serena MCP server \citep{Oraios2024Serena},
an open-source LSP-aggregator developed by Oraios AI that exposes
find-symbol, replace-symbol-body, and find-referencing-symbols
primitives to Claude Code over the Model Context Protocol.
\bibliography{refs}

\end{document}
